% ID: SEL-D0002-OA-E02
% ESTADO: validada. S003, 2026-11-07; tres pasadas del mismo autor, comprobación separada y revisión didáctica/visual.
\textbf{a) Variación en el intervalo permitido.}\par\nopagebreak[4]
Primero desarrollamos y derivamos:
\[
C(x)=8x^2-8x+3,\qquad C'(x)=16x-8=8(2x-1).
\]
La derivada es negativa si $x<1/2$ y positiva si $x>1/2$. En el dominio $[0,2]$, $C$ decrece en $[0,1/2]$ y crece en $[1/2,2]$ (los signos de la derivada se estudian en los interiores).

\textbf{b) Mínimo absoluto e interpretación.}\par\nopagebreak[4]
El cambio de signo de negativa a positiva da un mínimo en $x=1/2$. Comparamos también extremos:
\[
C(0)=3,\qquad C(1/2)=1,\qquad C(2)=19.
\]
Hay que producir \boxed{0.5\text{ millones de kg}=500\,000\text{ kg}}; el coste mínimo es \boxed{1}, en las unidades de coste del enunciado, que no especifica moneda ni escala monetaria. También se comprueba directamente: $2(2x-1)^2+1\geq1$, con igualdad únicamente en $x=1/2$.

\textbf{c) Esbozo.}\par\nopagebreak[4]
Es el arco de una parábola abierta hacia arriba, limitado a $0\leq x\leq2$, con vértice $(0.5,1)$ y eje $x=0.5$. Pasa por $(0,3)$, $(1,3)$ y $(2,19)$; $C''=16>0$. Los extremos del dominio están incluidos. No hay cortes con el eje horizontal, pues $C\geq1$.
\begin{center}\includegraphics[width=0.70\linewidth]{figuras/s003-coste.pdf}\end{center}
